\documentclass[aspectratio=169,xcolor=table]{beamer}
\usepackage[utf8]{inputenc}
\usepackage[T1]{fontenc}
\usepackage{lmodern}
\usepackage{csquotes}
\usepackage{xcolor}
\usepackage[portuguese]{babel}
\usepackage{hyperref}
\usepackage{tabularx}
\usepackage{amsmath}
\usepackage{amssymb}
\usepackage{tikz}
\usetikzlibrary{shapes.geometric, arrows.meta, positioning, fit, calc, mindmap, shadows}

% ------------------------------------------------
% Tema e Configurações do Beamer
% ------------------------------------------------
\usetheme{DCC}

% Ajuste de espaçamento entre itens
\setbeamertemplate{itemize items}[circle]
\setbeamertemplate{itemize subitem}[circle]
\setlength{\itemsep}{0.6em}
\setlength{\parskip}{0.4em}

\graphicspath{{imgs/}{./imgs/}}

\author[Luis, Antoniel, João]{%
  \textbf{Luis} \and \textbf{Antoniel} \and \textbf{João}
}
\title{Ambientes Virtuais de Aprendizagem}
\subtitle{Análise Crítica da Integração entre Tecnologia, Pedagogia e Prática Educacional}
\institute{Universidade Federal da Bahia \\ Instituto de Computação}
\date{13 de Dezembro de 2025}

\begin{document}

%-------------------------------------------------
%  SLIDE DE TÍTULO
%-------------------------------------------------
\begin{frame}[plain,noframenumbering]
    \titlepage
\end{frame}

%-------------------------------------------------
%  SLIDE DE AGENDA
%-------------------------------------------------
\begin{frame}{Agenda}
    \begin{table}
        \begin{tabularx}{\textwidth}{|l|X|}
            \hline
            \textbf{Tempo} & \textbf{Conteúdo} \\
            \hline
            0--5 min  & \textbf{Fundamentos dos AVAs} -- Conceitos, evolução histórica e funcionalidades (Luis Felipe) \\
            \hline
            5--10 min & \textbf{Abordagens Pedagógicas} -- Construtivismo vs Instrucionismo, TPACK e metodologias ativas (Antoniel) \\
            \hline
            10--15 min & \textbf{Análise Crítica} -- Benefícios, desafios, tendências e framework integrativo (João) \\
            \hline
        \end{tabularx}
    \end{table}
\end{frame}

\setlength{\parskip}{0.8em}

%=================================================
\section{Fundamentos dos AVAs}
%=================================================

\begin{frame}{O que são Ambientes Virtuais de Aprendizagem?}
    \begin{columns}[T]
        \begin{column}{0.55\textwidth}
            \begin{itemize}
                \item \textbf{Definição:} Plataformas digitais integradas para apoiar processos educacionais
                \item \textbf{Sinônimos:} LMS (Learning Management System), VLE (Virtual Learning Environment)
                \item \textbf{Função:} Gerenciar conteúdos, viabilizar interações e organizar atividades de ensino-aprendizagem
            \end{itemize}
        \end{column}
        \begin{column}{0.42\textwidth}
            \centering
            \begin{tikzpicture}[scale=0.7, transform shape]
                \node[draw, rounded corners, fill=blue!20, minimum width=2.5cm, minimum height=0.8cm] (ava) at (0,0) {\textbf{AVA}};
                \node[draw, rounded corners, fill=green!20, minimum width=2cm, minimum height=0.6cm] (cont) at (-2.5,1.5) {Conteúdo};
                \node[draw, rounded corners, fill=orange!20, minimum width=2cm, minimum height=0.6cm] (inter) at (2.5,1.5) {Interação};
                \node[draw, rounded corners, fill=red!20, minimum width=2cm, minimum height=0.6cm] (aval) at (-2.5,-1.5) {Avaliação};
                \node[draw, rounded corners, fill=purple!20, minimum width=2cm, minimum height=0.6cm] (colab) at (2.5,-1.5) {Colaboração};
                \draw[->, thick] (ava) -- (cont);
                \draw[->, thick] (ava) -- (inter);
                \draw[->, thick] (ava) -- (aval);
                \draw[->, thick] (ava) -- (colab);
            \end{tikzpicture}
        \end{column}
    \end{columns}
\end{frame}

\begin{frame}{Dimensões de Interação Pedagógica (Moore, 1989)}
    \begin{center}
        \begin{tikzpicture}[scale=0.85, transform shape,
            box/.style={draw, rounded corners, minimum width=3cm, minimum height=1cm, align=center, font=\small}]

            \node[box, fill=blue!30] (prof) at (0,2) {Professor};
            \node[box, fill=green!30] (aluno) at (-3,-1) {Aluno};
            \node[box, fill=orange!30] (conteudo) at (3,-1) {Conteúdo};

            \draw[<->, very thick, blue!60] (prof) -- node[left, font=\scriptsize] {Orientação e Feedback} (aluno);
            \draw[<->, very thick, green!60] (aluno) -- node[below, font=\scriptsize] {Exploração Autônoma} (conteudo);
            \draw[<->, very thick, orange!60] (prof) -- node[right, font=\scriptsize] {Design Instrucional} (conteudo);

            \node[box, fill=purple!20, dashed] (aluno2) at (-3,-3) {Aluno};
            \draw[<->, very thick, purple!60] (aluno) -- node[left, font=\scriptsize] {Colaboração} (aluno2);
        \end{tikzpicture}
    \end{center}
\end{frame}

\begin{frame}{Evolução Histórica da EaD}
    \begin{center}
        \begin{tikzpicture}[scale=0.75, transform shape]
            \draw[thick, ->, >=stealth] (0,0) -- (13,0) node[right] {\small Tempo};

            \foreach \x/\year/\tech/\col in {
                1/1728/Correio/gray!40,
                3.5/1920/Rádio/yellow!40,
                6/1950/TV/orange!40,
                8.5/1990/Internet/blue!40,
                11/2002/Moodle/green!40
            } {
                \draw[thick] (\x,0.1) -- (\x,-0.1);
                \node[above, font=\tiny] at (\x,0.2) {\year};
                \node[draw, rounded corners, fill=\col, below, font=\scriptsize, minimum width=1.2cm] at (\x,-0.5) {\tech};
            }
        \end{tikzpicture}
    \end{center}
    \vspace{0.3cm}
    \begin{itemize}
        \item \textbf{Séc. XIX:} Ensino por correspondência postal
        \item \textbf{Séc. XX:} Rádio e TV educativa expandem alcance
        \item \textbf{Anos 90:} Internet viabiliza sistemas integrados
        \item \textbf{2002:} Moodle -- filosofia construcionista social (Dougiamas)
        \item \textbf{2020-21:} Pandemia COVID-19 acelera adoção massiva
    \end{itemize}
\end{frame}

\begin{frame}{Funcionalidades dos AVAs Modernos}
    \begin{columns}[T]
        \begin{column}{0.48\textwidth}
            \textbf{Gestão de Conteúdo}
            \begin{itemize}
                \item Estruturação modular
                \item Materiais multimídia
                \item Hipertexto e recursos interativos
            \end{itemize}
            \vspace{0.5em}
            \textbf{Comunicação}
            \begin{itemize}
                \item Assíncrona: fóruns, e-mail
                \item Síncrona: chat, videoconferência
            \end{itemize}
        \end{column}
        \begin{column}{0.48\textwidth}
            \textbf{Avaliação}
            \begin{itemize}
                \item Quizzes automatizados
                \item Rubricas e portfólios
                \item Learning Analytics
            \end{itemize}
            \vspace{0.5em}
            \textbf{Colaboração}
            \begin{itemize}
                \item Wikis e glossários
                \item Projetos em grupo
                \item Avaliação entre pares
            \end{itemize}
        \end{column}
    \end{columns}
\end{frame}

%=================================================
\section{Abordagens Pedagógicas}
%=================================================

\begin{frame}{Instrucionismo vs Construtivismo}
    \begin{columns}[T]
        \begin{column}{0.48\textwidth}
            \textbf{\textcolor{red!70}{Instrucionismo}}
            \begin{itemize}
                \item AVA como repositório de conteúdo
                \item Professor transmite conhecimento
                \item Estudante é receptor passivo
                \item Foco em memorização
                \item ``Depósito de PDFs''
            \end{itemize}
        \end{column}
        \begin{column}{0.48\textwidth}
            \textbf{\textcolor{green!70!black}{Construtivismo}}
            \begin{itemize}
                \item AVA como espaço de mediação
                \item Professor é facilitador
                \item Estudante é protagonista
                \item Foco em construção ativa
                \item Colaboração e reflexão
            \end{itemize}
        \end{column}
    \end{columns}
    \vspace{0.5cm}
    \centering
    \textit{``A tecnologia per se não determina práticas pedagógicas'' (Jonassen, 1999)}
\end{frame}

\begin{frame}{Framework TPACK (Mishra \& Koehler, 2006)}
    \begin{columns}[T]
        \begin{column}{0.45\textwidth}
            \begin{center}
                \begin{tikzpicture}[scale=0.65, transform shape]
                    % Circles
                    \fill[blue!20, opacity=0.7] (0,0) circle (2cm);
                    \fill[red!20, opacity=0.7] (1.5,0) circle (2cm);
                    \fill[green!20, opacity=0.7] (0.75,1.3) circle (2cm);

                    % Labels
                    \node[font=\scriptsize] at (-1,0) {TK};
                    \node[font=\scriptsize] at (2.5,0) {PK};
                    \node[font=\scriptsize] at (0.75,2.5) {CK};
                    \node[font=\scriptsize\bfseries] at (0.75,0.5) {TPACK};
                \end{tikzpicture}
            \end{center}
        \end{column}
        \begin{column}{0.52\textwidth}
            \begin{itemize}
                \item \textbf{TK:} Conhecimento Tecnológico
                \item \textbf{PK:} Conhecimento Pedagógico
                \item \textbf{CK:} Conhecimento do Conteúdo
                \item \textbf{TPACK:} Integração contextualizada dos três domínios
            \end{itemize}
            \vspace{0.3cm}
            \small
            Formação docente deve transcender treinamentos puramente técnicos
        \end{column}
    \end{columns}
\end{frame}

\begin{frame}{Papéis Transformados}
    \begin{columns}[T]
        \begin{column}{0.48\textwidth}
            \textbf{Professor como Mediador}
            \begin{itemize}
                \item Seleciona e organiza recursos
                \item Projeta experiências de aprendizagem
                \item Modera discussões online
                \item Oferece feedback formativo
                \item Cria comunidade de aprendizagem
            \end{itemize}
        \end{column}
        \begin{column}{0.48\textwidth}
            \textbf{Estudante Autorregulado}
            \begin{itemize}
                \item Autonomia na gestão do tempo
                \item Proatividade na busca de recursos
                \item Competências colaborativas
                \item Habilidades metacognitivas
                \item Protagonismo na aprendizagem
            \end{itemize}
        \end{column}
    \end{columns}
\end{frame}

\begin{frame}{Metodologias Ativas em AVAs}
    \begin{center}
        \begin{tikzpicture}[scale=0.8, transform shape,
            metodbox/.style={draw, rounded corners, minimum width=2.8cm, minimum height=1cm, align=center, font=\scriptsize}]

            \node[metodbox, fill=blue!25] (pbl) at (0,0) {\textbf{PBL}\\Aprendizagem\\Baseada em\\Problemas};
            \node[metodbox, fill=green!25] (colab) at (4,0) {\textbf{Colaborativa}\\Wikis, Fóruns\\Projetos em\\Grupo};
            \node[metodbox, fill=orange!25] (flip) at (8,0) {\textbf{Sala Invertida}\\Conteúdo em\\Casa, Prática\\Presencial};
            \node[metodbox, fill=purple!25] (gam) at (4,-2.5) {\textbf{Gamificação}\\Pontos, Badges\\Rankings\\Desafios};
        \end{tikzpicture}
    \end{center}
    \vspace{0.3cm}
    \small
    \textbf{Santos Jr. (2023):} 90\% dos participantes reportaram aumento de motivação com gamificação em AVA
\end{frame}

%=================================================
\section{Análise Crítica}
%=================================================

\begin{frame}{Benefícios dos AVAs}
    \begin{columns}[T]
        \begin{column}{0.48\textwidth}
            \begin{itemize}
                \item \textbf{Democratização:} Rompe barreiras geográficas e temporais
                \item \textbf{Personalização:} Ritmos individuais de aprendizagem
                \item \textbf{Multimodalidade:} Múltiplas representações do conhecimento
            \end{itemize}
        \end{column}
        \begin{column}{0.48\textwidth}
            \begin{itemize}
                \item \textbf{Autonomia:} Desenvolvimento de autorregulação
                \item \textbf{Comunidades:} Interação entre contextos diversos
                \item \textbf{Analytics:} Dados para intervenções personalizadas
            \end{itemize}
        \end{column}
    \end{columns}
\end{frame}

\begin{frame}{Desafios e Limitações}
    \begin{columns}[T]
        \begin{column}{0.48\textwidth}
            \textbf{Aspectos Discentes}
            \begin{itemize}
                \item Isolamento e desmotivação
                \item Dificuldades de autorregulação
                \item Taxas de evasão elevadas
            \end{itemize}
            \vspace{0.5em}
            \textbf{Aspectos Docentes}
            \begin{itemize}
                \item Formação insuficiente
                \item Reprodução de práticas tradicionais
            \end{itemize}
        \end{column}
        \begin{column}{0.48\textwidth}
            \textbf{Aspectos Estruturais}
            \begin{itemize}
                \item Desigualdades digitais
                \item Design instrucional inadequado
                \item Vigilância e controle excessivo
            \end{itemize}
            \vspace{0.5em}
            \textbf{Aspectos Éticos}
            \begin{itemize}
                \item Privacidade de dados
                \item Exclusão digital
            \end{itemize}
        \end{column}
    \end{columns}
\end{frame}

\begin{frame}{Tendências Emergentes}
    \begin{center}
        \begin{tikzpicture}[scale=0.75, transform shape,
            trendbox/.style={draw, rounded corners, minimum width=2.5cm, minimum height=0.9cm, align=center, font=\scriptsize, drop shadow}]

            \node[trendbox, fill=blue!30] (hibrida) at (0,1.5) {\textbf{Educação}\\Híbrida};
            \node[trendbox, fill=green!30] (adapt) at (4,1.5) {\textbf{Aprendizagem}\\Adaptativa};
            \node[trendbox, fill=orange!30] (ia) at (8,1.5) {\textbf{Inteligência}\\Artificial};
            \node[trendbox, fill=purple!30] (mob) at (2,-0.5) {\textbf{Mobile}\\Learning};
            \node[trendbox, fill=red!30] (rea) at (6,-0.5) {\textbf{REA}\\Recursos Abertos};
            \node[trendbox, fill=cyan!30] (udl) at (4,-2.5) {\textbf{Design}\\Universal (UDL)};
        \end{tikzpicture}
    \end{center}
\end{frame}

\begin{frame}{Framework Conceitual Integrativo}
    \begin{center}
        \begin{tikzpicture}[scale=0.8, transform shape]
            % Three dimensions
            \node[draw, rounded corners, fill=blue!25, minimum width=3cm, minimum height=1.5cm, align=center] (tech) at (0,0) {\textbf{Dimensão}\\Tecnológica};
            \node[draw, rounded corners, fill=green!25, minimum width=3cm, minimum height=1.5cm, align=center] (ped) at (5,0) {\textbf{Dimensão}\\Pedagógica};
            \node[draw, rounded corners, fill=orange!25, minimum width=3cm, minimum height=1.5cm, align=center] (curr) at (2.5,-3) {\textbf{Dimensão}\\Curricular/Contextual};

            % Central node
            \node[draw, circle, fill=purple!30, minimum size=1.5cm, align=center, font=\small\bfseries] (efet) at (2.5,-1) {Efetividade\\Pedagógica};

            % Arrows
            \draw[->, thick] (tech) -- (efet);
            \draw[->, thick] (ped) -- (efet);
            \draw[->, thick] (curr) -- (efet);
        \end{tikzpicture}
    \end{center}
    \vspace{0.3cm}
    \small
    \centering
    A efetividade emerge da \textbf{integração coerente e contextualizada} das três dimensões
\end{frame}

\begin{frame}{Considerações Finais}
    \begin{itemize}
        \item AVAs são infraestruturas \textbf{sociotécnicas} centrais da educação contemporânea
        \item Efetividade não é determinada tecnologicamente, mas \textbf{socialmente construída}
        \item O desafio é \textbf{pedagógico, político e ético}, não apenas tecnológico
        \item Necessidade de articulação fundamentada entre \textbf{tecnologia, pedagogia e currículo}
    \end{itemize}
    \vspace{0.5cm}
    \centering
    \textit{``Como empregar AVAs de maneira que democratizem acesso ao conhecimento,\\ respeitem diversidades e promovam autonomia e pensamento crítico?''}
\end{frame}

%=================================================
\section{Referências}
%=================================================

\begin{frame}{Referências Principais}
    \tiny
    \begin{itemize}
        \item Moore, M. G. (1989). Three types of interaction. \textit{American Journal of Distance Education}, 3(2), 1-7.
        \item Garrison, D. R. (2016). \textit{E-Learning in the 21st Century}. Routledge.
        \item Mishra, P. \& Koehler, M. J. (2006). Technological Pedagogical Content Knowledge. \textit{Teachers College Record}, 108(6), 1017-1054.
        \item Jonassen, D. H. (1999). Designing constructivist learning environments. \textit{Instructional-design theories and models}. Erlbaum.
        \item Bates, A. W. (2015). \textit{Teaching in a Digital Age}. BCcampus.
        \item Dougiamas, M. (2003). Moodle: Using Learning Communities to Create an Open Source CMS. \textit{EdMedia Conference}.
        \item Hodges, C. et al. (2020). The Difference Between Emergency Remote Teaching and Online Learning. \textit{EDUCAUSE Review}.
    \end{itemize}
\end{frame}

%=================================================
\section{Encerramento}
%=================================================
\begin{frame}{Perguntas?}
    \begin{center}
        \Huge Obrigado!
        \vspace{1cm}

        \normalsize
        \textbf{Luis Felipe Sena} -- luis.sena@ufba.br\\
        \textbf{Antoniel Magalhães} -- antoniels@ufba.br\\
        \textbf{João Leahy} -- joao.leahy@ufba.br
    \end{center}
\end{frame}

\end{document}
